\documentclass[11pt]{article}
\usepackage[margin=0.75in]{geometry}
\usepackage{amsmath,amssymb,booktabs,array,xcolor,hyperref,listings}
\definecolor{constantred}{HTML}{DC2626}
\definecolor{codebg}{HTML}{F8FAFC}
\newcommand{\const}[1]{\textcolor{constantred}{#1}}
\newenvironment{narration}{\begin{quote}\itshape}{\end{quote}}
\lstset{
  basicstyle=\ttfamily\scriptsize,
  backgroundcolor=\color{codebg},
  frame=single,
  breaklines=true,
  columns=fullflexible,
  keepspaces=true
}
\setlength{\parindent}{0pt}
\setlength{\parskip}{0.5em}

\title{Three-Minute Technical Pitch Script}
\author{Shopping Copilot}
\date{}

\begin{document}
\maketitle

Read the italic narration. Equations, diagrams, tables, presenter cues, and animation directions \texttt{[like this]} guide the slide deck flow. Replace \texttt{[TEAM NAME]} before presenting.

\subsection*{Legend for On-Slide Colors and Cues}
\begin{itemize}
  \item \textbf{\const{Red Numbers} in Slide 1}: Fixed competition evaluation constants set by the organizers (not fitted by our team).
  \item \textbf{\const{Red Numbers} in Slides 3 \& 4}: Calibrated probabilistic hyperparameters (temperature, cutoffs, decay rates) governing our engine.
  \item \textbf{Red Dashed Outlines in Diagrams}: Optional/alternative components (e.g. LLM fallback) deliberately disabled in the submitted pipeline.
  \item \textbf{\texttt{[Animation / Visual Cue]}}: Direct cues for slide transitions, zooms, and spotlight animations.
\end{itemize}

\section*{Slide 1 - Outcome and problem (0:00-0:28)}

\textbf{[Screen Setup: Clean title slide displaying ``Shopping Copilot'' and key result badge: ``0.9744 TechnicalScore | 0 LLM Calls'']}

\begin{narration}
Hi everyone, we are [TEAM NAME]. We built Shopping Copilot---a deterministic, offline probabilistic
agent that finds a target product in a 50,000-item catalog in 10 turns or fewer. On the official public
leaderboard, we achieved a TechnicalScore of 0.9744 with zero LLM calls. Now, you might wonder if that's
just overfitting to the public set. To verify that, we tested on 2,800 target-disjoint test sessions
and scored 0.9562, plus an 800-session free-form opener stress test where we achieved 0.9345.
\end{narration}

\textbf{[Animation: Pop in the TechnicalScore formula with competition weights highlighted in RED]}

\[
\operatorname{TechnicalScore}
=\const{0.50}\,\operatorname{Hit@10}
+\const{0.30}\,\operatorname{MRR}
+\const{0.20}\,\operatorname{clip}\!\left(
\frac{\const{11}-\operatorname{MTTC}}{\const{10}},0,1\right)
\]

\textit{(Presenter Note: The red constants above are fixed competition weights defined by the organizer.)}

\textbf{[Animation: Slide in benchmark summary table with glowing highlight on the 0.9744, 0.9562, and 0.9345 scores]}

\begin{center}
\small
\begin{tabular}{lrrrrr}
\toprule
Reporting set & Sessions & Hit@10 & MRR & MTTC & TechnicalScore \\
\midrule
Official public & 200 & 1.0000 & 0.9942 & 2.19 & \textbf{0.9744} \\
Leakage-safe test & 2,800 & 0.9911 & 0.9783 & 2.64 & \textbf{0.9562} \\
Free-form-opener stress & 800 & 0.9725 & 0.9596 & 2.98 & \textbf{0.9345} \\
\bottomrule
\end{tabular}
\end{center}

\section*{Slide 2 - System architecture (0:28-0:58)}

\textbf{[Screen Setup: Picture A (High-Level Overview) appears first, then transitions into Picture B (Deep Technical Architecture)]}

\begin{narration}
When a customer message comes in, we run it through a fast, cheapest-first parsing cascade:
exact template matching, followed by ontology normalization.

You might wonder why we don't rely on standard tools like BM25, popularity priors, or an LLM fallback.
We actually built and tested all three. But our rigorous ablation studies proved that our two-level
probabilistic engine was already so precise that those extra components only added latency and noise
without improving accuracy. Disabling them gives us identical or higher accuracy, sub-50ms latency,
and zero LLM API cost.
\end{narration}

\textbf{[Animation: Flowchart animates left-to-right following the data path; Red dashed box on ``Disabled Components'' pulses with a ``DELIBERATELY ABLATED'' tag]}

\textbf{Picture A - High-level pipeline overview}

\begin{lstlisting}
flowchart LR
  Customer --> Cascaded_Parser_Templates_Ontology
  Cascaded_Parser_Templates_Ontology --> Session_State_Constraints_Decay
  Session_State_Constraints_Decay --> Level1_Category_Posterior
  Level1_Category_Posterior --> Level2_Item_Log_Posterior
  Level2_Item_Log_Posterior --> Expected_Utility_Depth_Policy
  Expected_Utility_Depth_Policy --> Ranked_ASINs_and_Question
\end{lstlisting}

\textbf{[Animation: Transition into Picture B --- Deep Technical Blueprint showing exact ML methods, mathematical gates, and probability bounds]}

\textbf{Picture B - Detailed technical architecture and mathematical flow}

\begin{lstlisting}
1. Deterministic NLP & State Tracking:
   Utterance -> Regex Template Match (0ms) -> If No -> Ontology Attribute Normalizer
   -> Session State Tracker (Constraints C_t, Age Decay gamma=0.9, Demote Factor=0.35)

2. Level 1 Bayesian Category Belief (1,115 Categories):
   Category Score s_c(x) = W_c(x)*coverage_c(x) + 3.0*1[quoted]*W_c(x)
   Softmax P(c|x) = softmax(s_c(x)/2.0 + 0.25*log(pi_c)) with Temperature T=2.0
   Prefix Mass Pruning (tau=0.85): sum P(c|x) >= 0.85 (Prunes 50k items -> ~250 candidates)

3. Level 2 Bounded Item Likelihood Fusion:
   Multi-Source Likelihoods: Exact Card (g=3.2), Soft-Card Jaccard (J>=0.34, g=1.5), Lexical Overlap
   Bounded Likelihood Floor: log L_r = log(max(0.02, exp(g_r*(s-1))))
   Item Log-Posterior: log P_t(i) ~ sum 0.9^(t-turn) [log L_main + log L_soft]
   Hard Exclusion Masking: Proven-inspected ASINs -> log P(i) = -infinity

4. Decision-Theoretic Recommendation Depth Policy:
   Dynamic Utility: k* = argmax [ sum (p_j/j) + (1 - sum p_j)*V ]
   Continuation Value: V = max(0, 0.75*d^s - 0.0667)
   Emit top k* recommendations (k* in 1..10, dynamically optimized per turn)
\end{lstlisting}

\textbf{Component Ablation and Pruning Analysis}

\begin{center}
\small
\begin{tabular}{p{3.5cm}p{4cm}p{4cm}p{3.5cm}}
\toprule
Candidate Component & Baseline / Role & Ablation Finding & Decision \& Production Benefit \\
\midrule
\textbf{BM25 / SQLite FTS5} & Official starter baseline & High noise across unrelated product categories (scored 0.1067) & \textbf{Disabled} (\texttt{gain=0.0}) --- replaced by Bayesian filter \\
\textbf{Item Popularity Prior} & E-commerce heuristic & Added $\Delta\text{Score} = +0.000000$ once soft-card evidence was active & \textbf{Disabled} (\texttt{weight=0.0}) --- removes catalogue bias \\
\textbf{Dense Embeddings (SVD/BLaIR)} & Vector similarity matching & Added memory/latency overhead with zero gain over ontology & \textbf{Disabled} --- keeps in-memory pipeline ultra-light \\
\textbf{Cloud LLM Fallback} & Edge-case prose parsing & Probabilistic engine scored 0.9562+; LLM added 50+ sec API latency & \textbf{Disabled} --- guarantees 0ms API latency, \$0 cost, 0 hallucinations \\
\bottomrule
\end{tabular}
\end{center}

\section*{Slide 3 - Two-level probabilistic ranking (0:58-1:42)}

\textbf{[Screen Setup: Split-view showing Level 1 (Category) on top and Level 2 (Item) on bottom. Highlight callout: ``Red numbers = Calibrated Probabilistic Hyperparameters'']}

\begin{narration}
For ranking, we do two-level probabilistic inference. At Level 1, we compute the category posterior
across all 1,115 categories using IDF-weighted token overlap and a catalog-share prior. Instead of picking
just one category and risking hard errors, we accumulate categories until we cover 85\% of the posterior mass.

At Level 2, we score individual items. Exact catalog cards, normalized attributes, lexical overlap, and
soft-card Jaccard all feed into an item log-posterior as bounded likelihoods. The red numbers you see
are our calibrated hyperparameters---like setting a likelihood floor of 0.02 so a noisy match never zeros
out the true target. Old evidence decays by a 0.9 factor each turn, and if a user changes their mind,
we demote the old constraint rather than completely wiping it.
\end{narration}

\textbf{[Animation: Zoom in on Level 1 equations; spotlight $\tau = 0.85$ and $T = 2.0$]}

\textbf{Level 1 - category posterior}

\[
s_c(x)=W_c(x)\,\operatorname{coverage}_c(x)
+\mathbb{1}[c\text{ quoted}]\,\const{3.0}\,W_c(x)
\]

\[
P(c\mid x)=\operatorname{softmax}\!\left(
\frac{s_c(x)}{T}
+\const{0.25}\log \pi_c\right),
\qquad
T=\const{2.0},
\qquad
\sum_{c\in\mathcal C_{\tau}}P(c\mid x)\geq \tau,
\quad \tau=\const{0.85}
\]

Here, $W_c$ is shared-token IDF mass and $\pi_c$ is the category's catalog share.

\textbf{[Animation: Pan down to Level 2 equations; highlight Likelihood floor $L_{\min} = 0.02$ and Age Weight $0.9^{\,t-\operatorname{turn}(e)}$]}

\textbf{Level 2 - bounded evidence and item posterior}

\[
\log L_r(i\mid e)=
\log\!\left[\max\!\left(
L_{\min},
\exp\{g_r(s_r(i,e)-1)\}
\right)\right],
\qquad L_{\min}=\const{0.02}
\]

\[
\log P_t(i\mid\mathcal D_t)\propto
\sum_{e\in\mathcal D_t}
\underbrace{\const{0.9}^{\,t-\operatorname{turn}(e)}}_{\text{age weight}}
\left[\log L_{\text{main}}(i\mid e)+\log L_{\text{soft}}(i\mid e)\right]
\]

\[
g_{\mathrm{exact}}=\const{3.2},\qquad
g_{\mathrm{soft}}=\const{1.5},\qquad
J_{\mathrm{soft}}\geq\const{0.34}
\]

If a returned item was definitely checked and the dialogue continued, then
$\log P_t(i)\leftarrow-\infty$. A demoted constraint receives an additional factor $\const{0.35}$.

\section*{Slide 4 - Recommendation depth (1:42-2:10)}

\textbf{[Screen Setup: Display Expected Utility formula $U(k)$ alongside an interactive curve plotting utility vs depth $k$]}

\begin{narration}
A big question is: how many products should we recommend at each turn? We treat this as an
expected-utility problem. For any depth k, our utility function balances the immediate reciprocal rank
against the expected continuation value of asking another question and getting cleaner evidence. We pick
the k-star that maximizes this trade-off. If we aren't getting new clues, the continuation value drops
and we widen our recommendations. And whenever the user rejects a suggested product, we know for sure it's
wrong, so we mask its log-posterior to negative infinity.
\end{narration}

\textbf{[Animation: On the utility curve, show the peak $k^*$ shifting dynamically from turn to turn; highlight negative infinity mask on rejected items]}

\[
U(k)=\sum_{j=1}^{k}\frac{p_j}{j}
+\left(1-\sum_{j=1}^{k}p_j\right)V,
\qquad
k^*=\underset{0\leq k\leq K}{\arg\max}\;U(k)
\]

\[
V=\max\!\left(0,\const{0.75}\,h-\const{0.0667}\right),
\qquad
h=d^{\,s},
\qquad
d=\begin{cases}
\const{0.8},&\text{templates are still matching},\\
\const{0.2},&\text{the deterministic parser is blind}.
\end{cases}
\]

\section*{Slide 5 - Optimization and evaluation discipline (2:10-2:35)}

\textbf{[Screen Setup: Clean optimization pipeline diagram with decision gate]}

\begin{narration}
To tune our hyperparameters, we hooked eight key constants into a Tree-structured Parzen Estimator
(TPE) offline. But to prevent overfitting to search noise, we set a strict rule: any change had to show
a statistically significant gain on a 95\% paired-bootstrap confidence interval. When recent proposals
didn't clear that bar, we kept our solid incumbent baseline. Best of all, running all 3,800 test sessions
through the official evaluator takes under 58 seconds total, with zero model calls.
\end{narration}

\textbf{[Animation: Highlight the 95\% CI formula; show green checkmark on solid incumbents and red X on noisy proposals that failed to beat zero]}

\[
x^*=\underset{x}{\arg\max}\;
\frac{\ell(x)}{g(x)},
\qquad
\text{adopt }x_j
\iff
\operatorname{CI}_{95\%}\!\left(\Delta\operatorname{TechnicalScore}_j\right)_{\mathrm{low}}>0
\]

Here, $\ell(x)=p(x\mid\text{good trials})$ and $g(x)=p(x\mid\text{remaining trials})$.

\begin{lstlisting}
Training split -> TPE proposal -> official-evaluator objective
               -> paired per-session bootstrap
               -> adopt only if the 95% CI clears zero
               -> otherwise keep the incumbent
\end{lstlisting}

\section*{Slide 6 - Live demo (2:35-3:00)}

\textbf{[Screen Setup: Full-screen browser window showing demo/index.html replay UI]}

\begin{narration}
Let's look at the live replay. In this templated browsing session, you can see how each requirement
sharpens the posterior distribution. Two checked items get eliminated, and our target hits rank 1 by turn
three. In this second case with a free-form opener, the first guess is wrong, but our agent recovers
immediately and surfaces the target at rank 1 on turn two. That sums up our whole approach: maintain
calibrated uncertainty, accumulate bounded evidence, and only recommend when expected utility says it's
worth it.
\end{narration}

\textbf{[Animation: Demo Replay 1 (Templated - browsing intent) --- Red bounding box zooms into pool size, excluded items, and highlights target reaching Rank 1 on Turn 3]}

\textbf{[Animation: Demo Replay 2 (Free-form - buying intent) --- Red spotlight highlights dynamic constraint update and target reaching Rank 1 on Turn 2]}

\section*{Accuracy notes - do not narrate unless asked}

\begin{itemize}
\item Do not call the submitted agent ``model-free.'' The accurate phrase is \textbf{deterministic,
offline probabilistic agent}.
\item The 12,000 sessions in the draft were an earlier fitting set, not a test set. The current split
is 8,400 training, 2,800 validation, and 2,800 test, disjoint by sample ID and target ASIN.
\item The 800-session stress test rewrites the opener; later turns retain official simulator wording.
It is not a fully natural multi-turn benchmark and is not an official competition score.
\item The latest recorded TPE proposals were rejected by the paired-bootstrap gate; do not imply
that they became submitted constants.
\item BM25, the item-popularity prior, and the optional LLM are disabled in the submitted path.
\end{itemize}

\end{document}
